%================================================================
\chapter{Introduction to \LaTeX{}}
\label{ch:latex}
%================================================================

\LaTeX{} is a document preparation system in which the author writes
plain text interleaved with markup commands, and a compiler transforms
that source into a typeset PDF. This contrasts with word processors
such as Microsoft Word or Google Docs, where formatting is applied
interactively in a what-you-see-is-what-you-get (WYSIWYG) editor.
The key property that makes \LaTeX{} the standard in science and
mathematics is its ability to render arbitrarily complex mathematical
notation with precision and consistency. Every document produced in
this program is a \LaTeX{} document.

%================================================================
\section{Your First \LaTeX{} Document}
\label{sec:first-latex-doc}
%================================================================

A \LaTeX{} source file is a plain text file with the extension
\texttt{.tex}. Just as \texttt{hello.py} was a plain text file that
Python interpreted, \texttt{hello.tex} is a plain text file that the
\LaTeX{} compiler processes into a PDF. The same text editor rules
apply (see Section~\ref{sec:plaintext-first}): use Notepad, TextEdit
in plain-text mode, or \texttt{gedit}; never a word processor.

\begin{infobox}{Exercise L-1: Your First \LaTeX{} Document}
Open a plain text editor (see Table~\ref{tab:text-editors}), type the
following exactly, and save the file as \texttt{hello.tex}:

\begin{lstlisting}[style=inlineStyle]
\documentclass{article}
\begin{document}

Hello, world. This is my first \LaTeX{} document.

Here is some mathematics: $e^{i\pi} + 1 = 0$.

\end{document}
\end{lstlisting}
\end{infobox}

Each line of the listing serves a specific purpose:

\begin{itemize}
    \item \verb|\documentclass{article}| declares the document class.
          For most reports and papers, \texttt{article} is the correct
          choice.
    \item \verb|\begin{document}| and \verb|\end{document}| delimit
          the body. Everything outside this pair is the
          \emph{preamble} (covered in
          Section~\ref{sec:latex-structure}).
    \item The text between the delimiters is the visible content.
    \item \verb|$e^{i\pi} + 1 = 0$| is inline math mode: text
          between dollar signs is typeset as mathematics.
\end{itemize}

\begin{table}[htbp]
\centering
\caption{Compiling a \LaTeX{} document on each platform}
\label{tab:latex-compile}
\footnotesize
\begin{tabular}{ll}
\toprule
\textbf{Platform} & \textbf{Command} \\
\midrule
Windows (Command Prompt or Git Bash) & \texttt{pdflatex hello.tex} \\
macOS (Terminal) & \texttt{pdflatex hello.tex} \\
Linux & \texttt{pdflatex hello.tex} or \texttt{latexmk -pdf hello.tex} \\
\bottomrule
\end{tabular}
\end{table}

Note that two compilation passes are often required for documents with
cross-references or a table of contents; \texttt{latexmk} handles
this automatically.

\vs
\noindent
Running \texttt{pdflatex hello.tex} produces several files. The
important one is \texttt{hello.pdf}, which can be opened in any PDF
viewer. The \texttt{.aux}, \texttt{.log}, and \texttt{.out} files are
auxiliary; they can be deleted safely.

%================================================================
\section{A More Complete Example}
\label{sec:latex-example2}
%================================================================

Exercise L-1 produced the minimal viable document. This section
introduces a more realistic document that demonstrates the features
most commonly needed in academic writing: a structured hierarchy of
sections, an automatically generated table of contents, in-text
citations linked to a bibliography, and the full range of mathematical
display modes. Every feature shown here is used in the documents
produced by this program.

\begin{infobox}{Exercise L-2: A Structured Document with Math and Bibliography}
Create two files in the same folder. The first is the main document,
\texttt{report.tex}. The second is the bibliography database,
\texttt{sample.bib}. Both files must be in the same directory for the
bibliography to resolve.

\vs
\noindent
\textbf{File 1: \texttt{report.tex}}

\begin{lstlisting}[style=inlineStyle]
\documentclass{article}

\usepackage{amsmath}   % extended math environments
\usepackage{booktabs}  % professional tables
\usepackage{hyperref}  % clickable links and cross-references

\title{A Short Report on Mathematical Notation}
\author{Your Name}
\date{\today}

\begin{document}

\maketitle
\tableofcontents
\newpage

\section{Introduction}

This report demonstrates document structure and mathematical
notation in \LaTeX{}. For background on literate programming,
see Knuth~\cite{knuth1984literate}. The \LaTeX{} system itself
is described in~\cite{lamport1994latex}.

\subsection{Motivation}

Scientific writing requires precise notation. \LaTeX{} provides
that precision through its mathematical typesetting engine.

\subsubsection{A Note on Style}

Use display equations for any expression that warrants visual
emphasis, and inline math for symbols embedded in prose.

\section{Mathematical Display Modes}

\subsection{Inline Math}

Inline math is enclosed in single dollar signs: $f(x) = x^2 + 1$.
It flows within the surrounding text without a line break.

\subsection{Display Math (unnumbered)}

The preferred form for unnumbered display equations is the
\texttt{\textbackslash[...\textbackslash]} environment:

\[
    \int_{-\infty}^{\infty} e^{-x^2}\,dx = \sqrt{\pi}
\]

The legacy \texttt{\$\$...\$\$} form produces the same visual
result but is deprecated in \LaTeX{} and should be avoided in
new documents.

\subsection{Numbered Equations}

The \texttt{equation} environment produces a numbered equation
that can be cross-referenced:

\begin{equation}
    \bar{x} = \frac{1}{n} \sum_{i=1}^{n} x_i
    \label{eq:mean}
\end{equation}

Equation~\ref{eq:mean} defines the sample mean. An approximation
for the standard deviation is given in~\cite{gurland1971simple}.

\subsection{Multi-line Aligned Equations}

The \texttt{align} environment aligns multiple equations at a
designated column, marked with \texttt{\&}:

\begin{align}
    \mu    &= \mathbb{E}[X]            \label{eq:mu}   \\
    \sigma^2 &= \mathbb{E}[(X-\mu)^2] \label{eq:var}
\end{align}

Equations~\ref{eq:mu} and~\ref{eq:var} define the mean and
variance of a random variable $X$.

\section{A Simple Table}

\begin{table}[htbp]
\centering
\caption{Comparison of equation display modes}
\label{tab:modes}
\begin{tabular}{lll}
\toprule
Mode & Syntax & Use case \\
\midrule
Inline     & \texttt{\$...\$}                    & Symbols in prose \\
Unnumbered & \texttt{\textbackslash[...\textbackslash]}  & Emphasis, no reference needed \\
Numbered   & \texttt{equation}                   & Cross-referenceable \\
Aligned    & \texttt{align}                      & Multi-line derivations \\
\bottomrule
\end{tabular}
\end{table}

\bibliographystyle{plain}
\bibliography{sample}

\end{document}
\end{lstlisting}

\vs
\noindent
\textbf{File 2: \texttt{sample.bib}}

\begin{lstlisting}[style=inlineStyle]
@article{knuth1984literate,
  author  = {Knuth, Donald E.},
  title   = {Literate Programming},
  journal = {The Computer Journal},
  year    = {1984},
  volume  = {27},
  number  = {2},
  pages   = {97--111}
}

@book{lamport1994latex,
  author    = {Lamport, Leslie},
  title     = {\LaTeX: A Document Preparation System},
  edition   = {2nd},
  publisher = {Addison-Wesley},
  year      = {1994}
}

@article{gurland1971simple,
  author  = {Gurland, John and Tripathi, Ram C.},
  title   = {A Simple Approximation for Unbiased Estimation
             of the Standard Deviation},
  journal = {The American Statistician},
  year    = {1971},
  volume  = {25},
  number  = {4},
  pages   = {30--32}
}
\end{lstlisting}
\end{infobox}

\begin{warningbox}{Compiling with a bibliography requires four steps}
A single \texttt{pdflatex} pass is not sufficient when the document
contains \verb|\cite{}| commands. The full sequence is:

\begin{verbatim}
pdflatex report.tex
bibtex report
pdflatex report.tex
pdflatex report.tex
\end{verbatim}

The first pass creates the \texttt{.aux} file listing all citation
keys. \texttt{bibtex} reads the \texttt{.aux} file and the
\texttt{.bib} database to produce a \texttt{.bbl} file containing
the formatted bibliography. The second and third \texttt{pdflatex}
passes incorporate the bibliography and resolve all
cross-references. Running \texttt{latexmk -pdf report.tex} performs
all necessary passes automatically.
\end{warningbox}

\begin{table}[htbp]
\centering
\caption{Equation display modes in \LaTeX{}}
\label{tab:equation-modes}
\begin{tabular}{p{2.8cm} p{4.2cm} p{2.5cm} p{4.5cm}}
\toprule
\textbf{Mode} & \textbf{Syntax} & \textbf{Output style} & \textbf{Notes} \\
\midrule
Inline & \verb|$...$| & Within text, compact & Use for symbols and short expressions in prose \\
Display (legacy) & \verb|$$...$$| & Centered, full size & Deprecated; avoid in new documents \\
Display (preferred) & \verb|\[...\]| & Centered, full size & Use for unnumbered display equations \\
Numbered & \verb|\begin{equation}| \verb|...\end{equation}| & Centered with number & Cross-referenceable via \verb|\label|/\verb|\ref| \\
Multi-line aligned & \verb|\begin{align}| \verb|...\end{align}| & Multiple lines, aligned at \verb|&| & Each line can carry its own \verb|\label| \\
\bottomrule
\end{tabular}
\end{table}

%================================================================
\section{Document Structure}
\label{sec:latex-structure}
%================================================================

Every \LaTeX{} source file has two parts: the \emph{preamble}
(everything before \verb|\begin{document}|) and the \emph{body}
(everything between \verb|\begin{document}| and
\verb|\end{document}|). The preamble declares the document class,
loads packages with \verb|\usepackage{...}|, and sets metadata. The
body contains the visible content, organized with sectioning commands.

\begin{lstlisting}[style=inlineStyle, caption={Minimal compilable document with preamble and sectioning}]
\documentclass{article}

% --- Preamble ---
\usepackage{amsmath}    % extended math
\usepackage{graphicx}   % figures
\usepackage{booktabs}   % professional tables
\usepackage{hyperref}   % clickable links

\title{My First Report}
\author{Alice Smith}
\date{\today}

% --- Body ---
\begin{document}

\maketitle

\section{Introduction}
This is the introduction.

\subsection{Background}
Some background material.

\subsubsection{Details}
More specific details.

\section{Methods}
Description of methods.

\end{document}
\end{lstlisting}

The \verb|\usepackage{...}| command loads extension packages. The
packages most relevant to this program are \texttt{amsmath} (extended
math), \texttt{graphicx} (figures), \texttt{booktabs} (tables), and
\texttt{hyperref} (clickable links). 

%================================================================
\section{Mathematics in \LaTeX{}: A Cheat Sheet}
\label{sec:latex-math}
%================================================================

\LaTeX{} distinguishes two math modes. Inline mode (\verb|$...$|)
typesets mathematics within a line of text. Display mode
(\verb|\[...\]| or the \texttt{equation} environment) centers
mathematics on its own line and assigns an equation number when using
\texttt{equation}. Use display mode for any expression that warrants
its own visual emphasis, and inline mode for symbols embedded in
prose.

%----------------------------------------------------------------
\subsection{Common Mathematical Notation}
%----------------------------------------------------------------

\begin{longtable}{p{3.2cm} p{4.5cm} p{3.0cm} p{3.5cm}}
\caption{Common \LaTeX{} mathematical notation} \label{tab:latex-math-cheatsheet} \\
\toprule
\textbf{Category} & \textbf{\LaTeX{} source} & \textbf{Rendered} & \textbf{Notes} \\
\midrule
\endfirsthead
\toprule
\textbf{Category} & \textbf{\LaTeX{} source} & \textbf{Rendered} & \textbf{Notes} \\
\midrule
\endhead
\midrule
\multicolumn{4}{r}{\emph{Continued on next page}} \\
\endfoot
\bottomrule
\endlastfoot

% --- Inline/display modes ---
\multicolumn{4}{l}{\textbf{Inline and display modes}} \\
\midrule
Inline math & \verb|$x + y = z$| & $x + y = z$ & Within a sentence \\
Display math & \verb|\[ x + y = z \]| & (centered block) & No number \\
Numbered equation & \verb|\begin{equation}...\end{equation}| & (with number) & Cross-referenceable \\
\midrule

% --- Basic arithmetic ---
\multicolumn{4}{l}{\textbf{Basic arithmetic}} \\
\midrule
Addition & \verb|a + b| & $a+b$ & \\
Subtraction & \verb|a - b| & $a-b$ & \\
Multiplication (dot) & \verb|a \cdot b| & $a \cdot b$ & Preferred in math \\
Multiplication (cross) & \verb|a \times b| & $a \times b$ & Vectors, cross product \\
Division & \verb|\frac{a}{b}| & $\frac{a}{b}$ & Always use \verb|\frac| in display \\
Superscript & \verb|x^{2}| & $x^{2}$ & Braces required for >1 char \\
Subscript & \verb|x_{i}| & $x_{i}$ & Braces required for >1 char \\
Both & \verb|x_{i}^{2}| & $x_{i}^{2}$ & \\
\midrule

% --- Common functions ---
\multicolumn{4}{l}{\textbf{Common functions}} \\
\midrule
Square root & \verb|\sqrt{x}| & $\sqrt{x}$ & \\
$n$th root & \verb|\sqrt[n]{x}| & $\sqrt[n]{x}$ & \\
Exponential & \verb|e^{x}| or \verb|\exp(x)| & $e^{x}$ & \verb|\exp| preferred in display \\
Natural log & \verb|\ln x| & $\ln x$ & \\
Log base $b$ & \verb|\log_{b} x| & $\log_{b} x$ & \\
Sine & \verb|\sin x| & $\sin x$ & Use \verb|\sin|, not \texttt{sin} \\
Cosine & \verb|\cos x| & $\cos x$ & \\
Absolute value & \verb+|x|+ or \verb|\lvert x \rvert| & $\lvert x \rvert$ & \verb|\lvert| scales with content \\
Floor & \verb|\lfloor x \rfloor| & $\lfloor x \rfloor$ & \\
Ceiling & \verb|\lceil x \rceil| & $\lceil x \rceil$ & \\
\midrule

% --- Summation and integration ---
\multicolumn{4}{l}{\textbf{Summation and integration}} \\
\midrule
Sum & \verb|\sum_{i=1}^{n} x_i| & $\sum_{i=1}^{n} x_i$ & Limits above/below in display \\
Product & \verb|\prod_{i=1}^{n} x_i| & $\prod_{i=1}^{n} x_i$ & \\
Integral & \verb|\int_{a}^{b} f(x)\,dx| & $\int_{a}^{b} f(x)\,dx$ & \verb|\,| adds thin space before $dx$ \\
Limit & \verb|\lim_{x \to 0} f(x)| & $\lim_{x \to 0} f(x)$ & \\
\midrule

% --- Greek letters ---
\multicolumn{4}{l}{\textbf{Greek letters (common in biostatistics and ML)}} \\
\midrule
alpha & \verb|\alpha| & $\alpha$ & Significance level, learning rate \\
beta & \verb|\beta| & $\beta$ & Regression coefficients \\
gamma & \verb|\gamma|, \verb|\Gamma| & $\gamma$, $\Gamma$ & Gamma distribution \\
delta & \verb|\delta|, \verb|\Delta| & $\delta$, $\Delta$ & Change, Dirac delta \\
epsilon & \verb|\epsilon|, \verb|\varepsilon| & $\epsilon$, $\varepsilon$ & Small quantity \\
theta & \verb|\theta|, \verb|\Theta| & $\theta$, $\Theta$ & Parameters, angle \\
lambda & \verb|\lambda|, \verb|\Lambda| & $\lambda$, $\Lambda$ & Poisson rate, eigenvalue \\
mu & \verb|\mu| & $\mu$ & Mean \\
nu & \verb|\nu| & $\nu$ & Degrees of freedom \\
pi & \verb|\pi|, \verb|\Pi| & $\pi$, $\Pi$ & 3.14\ldots, product \\
rho & \verb|\rho| & $\rho$ & Correlation \\
sigma & \verb|\sigma|, \verb|\Sigma| & $\sigma$, $\Sigma$ & Std dev, summation \\
tau & \verb|\tau| & $\tau$ & Time constant \\
phi & \verb|\phi|, \verb|\Phi| & $\phi$, $\Phi$ & Normal CDF, angle \\
chi & \verb|\chi| & $\chi$ & Chi-squared test \\
omega & \verb|\omega|, \verb|\Omega| & $\omega$, $\Omega$ & Frequency, sample space \\
\midrule

% --- Relations and logic ---
\multicolumn{4}{l}{\textbf{Relations and logic}} \\
\midrule
Not equal & \verb|\neq| & $\neq$ & \\
Less/greater or equal & \verb|\leq|, \verb|\geq| & $\leq$, $\geq$ & \\
Approximately & \verb|\approx| & $\approx$ & \\
Proportional to & \verb|\propto| & $\propto$ & \\
Distributed as & \verb|\sim| & $\sim$ & \\
For all & \verb|\forall| & $\forall$ & \\
There exists & \verb|\exists| & $\exists$ & \\
Implies & \verb|\Rightarrow| & $\Rightarrow$ & \\
If and only if & \verb|\Leftrightarrow| & $\Leftrightarrow$ & \\
Element of & \verb|\in| & $\in$ & \\
Subset & \verb|\subset|, \verb|\subseteq| & $\subset$, $\subseteq$ & \\
Union, intersection & \verb|\cup|, \verb|\cap| & $\cup$, $\cap$ & \\
\midrule

% --- Number sets ---
\multicolumn{4}{l}{\textbf{Number sets}} \\
\midrule
Real numbers & \verb|\mathbb{R}| & $\mathbb{R}$ & Requires \texttt{amsfonts} or \texttt{amssymb} \\
Natural numbers & \verb|\mathbb{N}| & $\mathbb{N}$ & \\
Integers & \verb|\mathbb{Z}| & $\mathbb{Z}$ & \\
Rationals & \verb|\mathbb{Q}| & $\mathbb{Q}$ & \\
Complex & \verb|\mathbb{C}| & $\mathbb{C}$ & \\
\midrule

% --- Vectors and matrices ---
\multicolumn{4}{l}{\textbf{Vectors and matrices}} \\
\midrule
Bold vector & \verb|\mathbf{x}| & $\mathbf{x}$ & Preferred in statistics \\
Vector with arrow & \verb|\vec{x}| & $\vec{x}$ & Physics convention \\
Transpose & \verb|A^{\top}| or \verb|A^T| & $A^{\top}$ & \\
$2\times 2$ matrix & \verb|\begin{bmatrix} a & b \\ c & d \end{bmatrix}| & $\begin{bmatrix} a & b \\ c & d \end{bmatrix}$ & Requires \texttt{amsmath} \\
Norm & \verb+\|x\|+ & $\|x\|$ & \\
Inner product & \verb|\langle x, y \rangle| & $\langle x, y \rangle$ & \\

\end{longtable}

%================================================================
\section{Figures}
\label{sec:latex-figures}
%================================================================

Figures are included using the \texttt{figure} float environment and
the \verb|\includegraphics| command from the \texttt{graphicx}
package. The \texttt{pdflatex} compiler accepts PDF, PNG, and JPG image formats.

\begin{lstlisting}[style=inlineStyle, caption={Including a figure in \LaTeX{}}]
\begin{figure}[htbp]
    \centering
    \includegraphics[width=0.7\textwidth]{images/my-plot.png}
    \caption{Description of the figure.}
    \label{fig:my-plot}
\end{figure}
\end{lstlisting}

The \texttt{[htbp]} option suggests placement preferences to
\LaTeX{}: here, top, bottom, or a dedicated page. The
\verb|\caption| and \verb|\label| commands provide a numbered caption
and a cross-reference target, respectively.

%================================================================
\section{Tables}
\label{sec:latex-tables}
%================================================================

Tables are built with the \texttt{table} float and the
\texttt{tabular} environment. The \texttt{booktabs} package 
provides clean horizontal rules.
Vertical rules (\texttt{|}) are omitted by convention in professional
documents; \texttt{booktabs} horizontal rules are sufficient.

\begin{lstlisting}[style=inlineStyle, caption={A minimal table with booktabs}]
\begin{table}[htbp]
\centering
\caption{Sample measurements}
\label{tab:sample}
\begin{tabular}{lrr}
\toprule
Subject & Height (cm) & Weight (kg) \\
\midrule
Alice   & 165         & 60          \\
Bob     & 178         & 75          \\
Carol   & 172         & 68          \\
\bottomrule
\end{tabular}
\end{table}
\end{lstlisting}

%================================================================
\section{Bibliography}
\label{sec:latex-bibliography}
%================================================================

\LaTeX{} manages bibliographies via an external file (\texttt{.bib})
and a style command. The \verb|\cite{key}| command inserts a citation
at any point in the text; the \verb|\bibliography{filename}| and
\verb|\bibliographystyle{plain}| commands generate the reference list.

For more sophisticated bibliography management, the \texttt{biblatex}
package with the \texttt{biber} backend is preferred in modern
documents. This repository's shared preamble has bibliography support
commented out; students writing independent reports should uncomment
and configure it.

%================================================================
\section{LaTeX Workshop in Visual Studio Code}
\label{sec:latex-workshop}
%================================================================

The LaTeX Workshop extension for VSC (referenced in
Section~\ref{sec:extensions}) provides a complete integrated workflow:
syntax highlighting, live PDF preview, one-click build, and error
navigation. Once installed, opening any \texttt{.tex} file activates
the extension automatically. The build shortcut is
\texttt{Ctrl+Alt+B} (Windows/Linux) or \texttt{Cmd+Option+B}
(macOS), which runs \texttt{pdflatex} and displays the result in the
built-in PDF viewer. Students do not need to use the command line for
compilation after this is set up; the command-line approach in
Section~\ref{sec:first-latex-doc} is presented first so the
underlying process is understood.

%================================================================
\section{Summary}
\label{sec:latex-summary}
%================================================================

This chapter introduced \LaTeX{}, the document preparation system
used throughout this program and widely adopted as the standard in
science and mathematics:

\begin{itemize}
    \item \textbf{What \LaTeX{} is.} A plain-text markup language
          compiled into typeset PDFs, in contrast to WYSIWYG word
          processors.
    \item \textbf{The source-to-PDF workflow.} Writing a
          \texttt{.tex} file in a text editor, compiling with
          \texttt{pdflatex} or \texttt{latexmk}, and viewing the
          resulting PDF.
    \item \textbf{Document structure.} The preamble (document class,
          packages, metadata) and the body (content organized with
          \verb|\section|, \verb|\subsection|, and
          \verb|\subsubsection|).
    \item \textbf{Inline and display math.} Dollar signs for inline
          mathematics and \verb|\[...\]| or the \texttt{equation}
          environment for display mathematics.
    \item \textbf{Equation modes.} Five modes exist: inline
          (\verb|$...$|), legacy display (\verb|$$...$$|, deprecated),
          preferred unnumbered display (\verb|\[...\]|), numbered
          (\texttt{equation}), and multi-line aligned (\texttt{align}).
          Use \verb|\[...\]| for unnumbered and \texttt{equation} for
          cross-referenceable display equations.
    \item \textbf{The math cheat sheet.} A comprehensive reference
          table (Table~\ref{tab:latex-math-cheatsheet}) covering
          arithmetic, functions, summation, integration, Greek
          letters, relations, number sets, and vectors/matrices.
    \item \textbf{Figures and tables.} Float environments with
          \verb|\includegraphics|, \texttt{tabular}, and
          \texttt{booktabs} for professional formatting.
    \item \textbf{Bibliography workflow.} A \texttt{.bib} file stores
          reference entries. The \verb|\cite{key}| command inserts a
          citation; \verb|\bibliography{filename}| and
          \verb|\bibliographystyle{plain}| generate the reference list.
          Documents with citations require a four-step compilation
          sequence: two \texttt{pdflatex} passes enclosing one
          \texttt{bibtex} pass, followed by a final \texttt{pdflatex}.
          \texttt{latexmk} handles this automatically.
    \item \textbf{LaTeX Workshop.} The VSC extension that provides
          syntax highlighting, live preview, and one-click build,
          eliminating the need for command-line compilation in daily
          use.
\end{itemize}
